% Source: https://artofproblemsolving.com/community/c3h1942324_perfect_square
% Stable identity: aops-topic-1942324-q01
Write $x^3-x^2=x^2(x-1)$. For $x\ge1$, $x^2$ is a nonzero perfect square, so $x^2(x-1)$ is a perfect square exactly when $x-1$ is a perfect square. (If $x^2m=n^2$ with $x\ne0$, then $m=\left(\dfrac {n}{x}\right)^2$ is the square of a rational number and an integer, hence of an integer.)

With $1\le x\le100$ we need $x-1\in\{0,1,4,9,\dots,81\}$, the squares $0^2,\dots,9^2$ (the next, $100$, gives $x=101$). So there are $\boxed{10}$ such integers ($x=1,2,5,10,17,26,37,50,65,82$).
